\documentclass[11pt, a4paper]{article}

\usepackage{amsmath, amssymb}
\usepackage{booktabs}
\usepackage{xcolor}
\usepackage{geometry}
\usepackage{hyperref}

\geometry{margin=1in}

\title{\textbf{Laketex Theming Engine \& Custom Themes Guide}}
\author{Laketex Design System Team}
\date{October 2026}

\begin{document}

\maketitle

\begin{abstract}
This guide provides complete technical specifications for the Laketex key-value theming engine. It explains how to author custom JSON and YAML theme files, configure color palettes, place them into the \texttt{themes/} folder alongside the executable, and activate them inside the user interface.
\end{abstract}

\section{Introduction to the Theming Architecture}

Laketex decouples user interface aesthetics from code compilation through a flexible, key-value configuration system ([`ThemeConfig`]). All interface colors—including panels, buttons, editor text, selection highlights, carets, borders, and compiler status indicators—are parameterized through standard hexadecimal strings.

\section{Theme File Discovery \& The \texttt{themes/} Folder}

When Laketex starts, it automatically traverses three discovery locations for theme files:
\begin{enumerate}
    \item \textbf{Application \texttt{themes/} Directory}: \texttt{./themes/} located alongside \texttt{laketex.exe}.
    \item \textbf{Working Directory \texttt{themes/} Directory}: \texttt{./themes/} in the current document root.
    \item \textbf{Root Executable Directory}: Loose \texttt{.json}, \texttt{.yaml}, or \texttt{.yml} files placed directly alongside \texttt{laketex.exe}.
\end{enumerate}

Any valid theme file discovered in these paths is dynamically loaded into memory and populated in the top toolbar's theme selector dropdown.

\section{Key-Value Color Specification}

Theme files support simple key-value definitions. You do not need to specify every key; unspecified keys gracefully fall back to sensible defaults.

\begin{table}[htbp]
\centering
\begin{tabular}{lll}
\toprule
\textbf{Token Key} & \textbf{Accepted Synonyms} & \textbf{Component Styled} \\
\midrule
\texttt{background} & \texttt{bg}, \texttt{window\_fill} & Outer application window \\
\texttt{panel\_background} & \texttt{panel\_bg}, \texttt{surface} & Sidebar, status bar, and console drawer \\
\texttt{editor\_background} & \texttt{editor\_bg}, \texttt{code\_bg} & Central code editor text buffer \\
\texttt{text} & \texttt{foreground}, \texttt{fg} & Primary readable text and labels \\
\texttt{text\_muted} & \texttt{comment}, \texttt{muted} & Comments, line numbers, subtle hints \\
\texttt{selection} & \texttt{selection\_bg} & Active editor highlight rectangle \\
\texttt{accent} & \texttt{primary}, \texttt{cursor} & Editor caret, active links, primary buttons \\
\texttt{border} & \texttt{stroke}, \texttt{separator} & Panel dividers and group frames \\
\texttt{error} & \texttt{err} & Diagnostic error indicators (\#ff5555) \\
\texttt{warning} & \texttt{warn} & Diagnostic warning indicators (\#ffb86c) \\
\texttt{success} & \texttt{ok} & Compile success indicators (\#50fa7b) \\
\midrule
\texttt{syntax\_command} & \texttt{command} & Control commands (\textbackslash begin, \textbackslash section, \textbackslash textbf) \\
\texttt{syntax\_argument} & \texttt{argument} & Command and environment arguments (\{itemize\}, \{article\}) \\
\texttt{syntax\_package} & \texttt{syntax\_package\_name} & Imported packages in \textbackslash usepackage\{amsmath, graphicx\} \\
\texttt{syntax\_tag\_property} & \texttt{syntax\_property} & Tag properties (leftmargin, label in \textbackslash begin\{itemize\}[...]) \\
\texttt{syntax\_property\_number} & \texttt{syntax\_number} & Integer and dimension values of tag properties (0.15in, 12pt) \\
\texttt{syntax\_property\_boolean} & \texttt{syntax\_boolean} & Boolean values of tag properties (true, false) \\
\texttt{syntax\_property\_value} & \texttt{syntax\_value} & General values of tag properties \\
\texttt{syntax\_math} & \texttt{math} & Inline math blocks (\$ ... \$) \\
\texttt{syntax\_math\_block} & \texttt{math\_block} & Display math blocks (\$\$ ... \$\$) \\
\texttt{syntax\_brace} & \texttt{brace} & Structural curly braces (\{ and \}) \\
\texttt{syntax\_bracket} & \texttt{bracket} & Optional argument brackets ([ and ]) \\
\texttt{syntax\_special} & \texttt{special} & Special control characters (\&, \_, \^{}, \textasciitilde, \#) \\
\texttt{syntax\_comment} & \texttt{comment} & Comments (\% ...) \\
\texttt{syntax\_text} & \texttt{text} & Regular prose text \\
\texttt{bracket\_level\_0..3} & -- & Depth-aware rainbow bracket levels 0 through 3 \\
\bottomrule
\end{tabular}
\caption{Complete reference of recognized theme color keys.}
\label{tab:theme_keys}
\end{table}

\section{Format Specifications \& Examples}

Themes can be authored in either standard JSON or clean YAML syntax.

\subsection{JSON Theme Template (\texttt{themes/my\_theme.json})}
\begin{verbatim}
{
  "name": "Nordic Slate",
  "dark_mode": true,
  "colors": {
    "background": "#242933",
    "panel_background": "#2e3440",
    "editor_background": "#1e222a",
    "text": "#eceff4",
    "text_muted": "#7b88a1",
    "accent": "#88c0d0",
    "selection": "#434c5e",
    "border": "#3b4252",
    "error": "#bf616a",
    "warning": "#ebcb8b",
    "success": "#a3be8c"
  }
}
\end{verbatim}

\subsection{YAML Theme Template (\texttt{themes/my\_theme.yaml})}
\begin{verbatim}
name: Sunset Warmth
dark_mode: true
colors:
  background: "#1f1b24"
  panel_background: "#26212d"
  editor_background: "#18151c"
  text: "#f5e6c8"
  text_muted: "#9e8b7d"
  accent: "#ff9e64"
  selection: "#443850"
  border: "#3a3044"
  error: "#f7768e"
  warning: "#e0af68"
  success: "#9ece6a"
\end{verbatim}

\section{Built-in JetBrains Dracula Theme}

For users accustomed to IntelliJ, PyCharm, or CLion, Laketex includes an official **Dracula (JetBrains)** theme built directly into the program:
\begin{itemize}
    \item Editor Canvas: \texttt{\#282a36}
    \item Window Panels: \texttt{\#21222c}
    \item Selection Background: \texttt{\#44475a}
    \item Purple Accent: \texttt{\#bd93f9}
    \item Foreground Text: \texttt{\#f8f8f2}
    \item Comments: \texttt{\#6272a4}
\end{itemize}

\section{Testing \& Hot-Reloading Themes}

When you place a new theme file into the \texttt{themes/} directory, you can immediately select it from the top toolbar dropdown. If you modify a custom theme file on disk, simply re-select it in the dropdown to instantly refresh the active interface.

\end{document}
